% from ega3-1-fr.tex
\oldpage[III]{125}

\[
\xymatrix{
H^n(X,\mathcal F)\ar[r]\ar[dr]&
H^n(\widehat X,\widehat{\mathcal F})
=H^n(\widehat X,i^*(\mathcal F))\ar[d]\\
&H^n(X,\mathcal F_k)=H^n(X_k,i_k^*(\mathcal F))
}
\]
d'où on déduit aussitôt la commutativité de
(\hyperref[III.4.1.4.3-fr]{4.1.4.3}).
\end{env}

\begin{theorem}[4.1.5]
\label{III.4.1.5-fr}
Soient $f:X\to Y$ un morphisme propre de préschémas noethériens, $Y'$ une
partie fermée de $Y$, $X'=f^{-1}(Y')$. Alors, pour tout
$\mathcal O_X$-Module cohérent $\mathcal F$,
$R^n\widehat f_*(\widehat{\mathcal F})$ est un
$\mathcal O_{\widehat X}$-Module cohérent et les homomorphismes
$\varphi_n$, $\psi_n$ et $\rho_n$ du diagramme
(\hyperref[III.4.1.4.2-fr]{4.1.4.2}) sont des isomorphismes topologiques.
\end{theorem}

\begin{proof}
Il suffira évidemment de prouver que $\varphi_n$ et $\psi_n$ sont des
isomorphismes~; comme $R^nf_*(\mathcal F)$ est cohérent
(\hyperref[III.3.2.1-fr]{3.2.1}), il en résultera que
$\widehat{R^nf_*(\mathcal F)}$ est cohérent (I, 10.8.8) et la bicontinuité
de $\varphi_n$, $\psi_n$ et $\rho_n$ est alors automatique (I, 10.11.6).
\end{proof}

\begin{env}[4.1.6]
\label{III.4.1.6-fr}
\emph{Remarques.}
\begin{enumerate}
\item[(i)] Si on pose
$\widehat{\mathcal F}_k=
\widehat{\mathcal F}/\widehat{\mathcal K}^{k+1}
\widehat{\mathcal F}$, il est immédiat que
$\mathcal F_k=i_*(\widehat{\mathcal F}_k)$, et l'homomorphisme canonique
(\hyperref[III.4.1.3.6-fr]{4.1.3.6}) n'est autre que l'homomorphisme déjà
défini en (\hyperref[III.3.4.2.2-fr]{3.4.2.2})
\[
R^n\widehat f_*(\widehat{\mathcal F})
\longrightarrow\varprojlim_kR^n\widehat f_*(\widehat{\mathcal F}_k)~;
\tag{4.1.6.1}\label{III.4.1.6.1-fr}
\]
par suite le fait que $\psi_n$ soit un isomorphisme est un cas particulier
de (\hyperref[III.3.4.3-fr]{3.4.3}). Mais nous en donnerons ci-dessous une
démonstration directe, évitant les considérations délicates sur les limites
projectives de suites spectrales (0, 13.7), sur lesquelles repose le
théorème général (\hyperref[III.3.4.3-fr]{3.4.3}).

\item[(ii)] Compte tenu du fait que les $\psi_n$ sont des isomorphismes, il
est équivalent de dire que les $\varphi_n$ ou les
$\rho_n=\psi_n^{-1}\circ\varphi_n$ sont des isomorphismes. Le th.
(\hyperref[III.4.1.5-fr]{4.1.5}) exprime entre autres que la formation des
$R^nf_*$ commute à la complétion et pourra s'appeler le premier théorème de
comparaison de la théorie «~algébrique~» à la théorie «~formelle~».
\end{enumerate}
Nous commencerons par établir la forme affine de
(\hyperref[III.4.1.5-fr]{4.1.5})~:
\end{env}

\begin{corollary}[4.1.7]
\label{III.4.1.7-fr}
Les hypothèses étant celles de (\hyperref[III.4.1.5-fr]{4.1.5}), supposons
en outre $Y=\operatorname{Spec}(A)$, où $A$ est noethérien, et
$\mathcal J=\widetilde{\mathfrak J}$, où $\mathfrak J$ est un idéal de
$A$, de sorte que
$\mathcal F_k=\mathcal F/\mathfrak J^{k+1}\mathcal F$.
L'homomorphisme canonique
\[
\varphi_n:\widehat{H^n(X,\mathcal F)}
\longrightarrow\varprojlim_kH^n(X,\mathcal F_k)
\tag{4.1.7.1}\label{III.4.1.7.1-fr}
\]
(où le premier membre est le séparé complété de $H^n(X,\mathcal F)$ pour
la topologie $\mathfrak J$-préadique) est un isomorphisme. Le système
projectif $(H^n(X,\mathcal F_k))_{k\geq0}$ vérifie pour tout $n$ la
condition (ML), et l'homomorphisme canonique
\[
\psi_n:H^n(\widehat X,\widehat{\mathcal F})
\longrightarrow\varprojlim_kH^n(X,\mathcal F_k)
\tag{4.1.7.2}\label{III.4.1.7.2-fr}
\]

